\documentclass[10pt]{article}

\usepackage[utf8]{inputenc}

\usepackage[T1]{fontenc}

\usepackage{amsmath}

\usepackage{amsfonts}

\usepackage{amssymb}

\usepackage[version=4]{mhchem}

\usepackage{stmaryrd}

\usepackage{graphicx}

\usepackage[export]{adjustbox}

\graphicspath{ {./images/} }

\usepackage{mathrsfs}

\usepackage{CJKutf8}



\begin{document}

в задаче выбора решения $d$ из пространства решений $D$, то наибольший ожидаемый риск



\[

\sup _{\theta \in \Theta} R\left(\theta, \delta^{*}\right)

\]



будет настолько малым, насколько это возможно. Принцип М. с. п. не всегда приводит к разумным выводам



\includegraphics[max width=\textwidth, center]{2023_07_08_f09922783b3f583efe9cg-1(5)}

(см. рис.); в данном случае следует ориентироваться на процедуру $\delta_{1}$, а не на $\delta_{2}$, хотя $\sup _{\theta \in \Theta} R\left(\theta, \delta_{1}\right)>\sup _{\theta \in \Theta} R\left(\theta, \delta_{2}\right)$



Понятие М. с. П. является полезным в задачах принятия статистич. решений в условиях отсутствия априорной информации относительно параметра $\theta$.



Лит.: [1] Л е м а н Э., Проверка статистических гипотез, пер. с англ., 2 изд., М., 1979; [2] 3 а к с ШІ., Теория статистиче ских выводов, пер. с англ., М., 1975. М. С. Никулил.



МИНИМАЛЬНАЯ ДОСТАТОЧНАЯ СТАТИСТИКА, не обходимая дост а тотная с т а т ис т и к а,- статистика $X$, являющаяся достаточной статистикой для семейства распределений $\mathscr{P}=\left\{P_{\theta}\right.$, $\theta \in \Theta\}$ и такая, что для любой другой достаточной статистики $Y$ имеет место $X=g(Y)$, где $g$ - нек-рая измеримая функция. Достаточная статистика минимальна тогда и только тогда, когда порождаемая ею достаточная $\sigma$-алгебра минимальна, т. е. содержится в любой другой достаточной б-алгебре.



Используется также понятие $\mathscr{P}-$ и и и м а л ь н о й



\includegraphics[max width=\textwidth, center]{2023_07_08_f09922783b3f583efe9cg-1}

р ы). Достаточная $\sigma$-алгебра $\mathscr{B}_{0}$ (и соответствующая ста-



\includegraphics[max width=\textwidth, center]{2023_07_08_f09922783b3f583efe9cg-1(4)}

пополнении $\overline{\mathscr{B}}$ любой достаточной б-алгебры $\mathscr{B}$ относительно семейства распределений $\mathscr{P}$. Если семей-



\includegraphics[max width=\textwidth, center]{2023_07_08_f09922783b3f583efe9cg-1(7)}

гебра $\mathscr{B}_{0}$, порождаемая семейством плотностей



\[

\left\{p_{\theta}(\omega)=\frac{d P}{d \mu}(\omega) ; \theta \in \Theta\right\}

\]



является достаточной и $\mathscr{P}$ минимальной.



Общий пример М. д. с. дает направляющая статистика $T=\left(T_{1}, \ldots, T_{n}\right)$ әкспоненциального семейства



\[

p_{\theta}(\omega)=C(\theta) \exp \sum_{j} Q_{j}(\theta) T_{j}(\omega)

\]



в предположении тотности канонич. параметризации этого семейства.



Лит.: [1] Б а p p a Ж.-Р., Основные понятия математической статистики, пер. с франц., М., 1974; [2] III м е т т е p е p J.,



A. С.'Холева.



МИНИМАЛЬНАЯ МОДЕЛЬ - алгебраическое многообразие с условием минимальности относительно существования бирационалыных морфизмов на неособые многообразия. Точнее, пусть $\boldsymbol{B}$ - класс всех бирационально әквивалентных неособых проективных многообразий над алгебраически замкнутым полем $k$, поля функций к-рых изоморфны заданному конечно порожденному расширению $K$ над $k$. Многообразия из класса $B$ наз. П р о е т и в ны м и мо д л я м и әтого класса, или проективными моделями поля $K / k$. Многообразие



\includegraphics[max width=\textwidth, center]{2023_07_08_f09922783b3f583efe9cg-1(6)}

м о е л ь ю, если всякий бирационалыный морфизм $f: X \rightarrow X^{\prime}$, где $X^{\prime} \in B$, является изоморфизмом. Иначе говоря, относительно М. м.-- это минимальные элементы в $B$ относитслыно частичного порядка, определяемого следующим отноцением доминирования: $X_{1}$ доминирует $X_{2}$, если существует бирациональный морфизм $h: X_{1} \rightarrow X_{2}$. Если относительно М. М. елинственна в $B$, то она наз. м и п ма л ь о й моде лью. В каждом классе бирационально эквивалентных кривых сугествует единственная (с точностью до изоморфизма) неособая проективная кривая. Так что всякая неособая проективная кривая является М. М. В общем случае, если класс $B$ не пуст, то в нем существует хотя бы одиа относительно М. м. Непустота класса $B$ известна (благодаря теоремам о разрешении особенностей) для многообразий любой размерности в характеристике 0 и для многообразий размерности $n \leqslant 3$ в характеристике $p>5$.



Основные результаты о М. м. алгебраич. поверхностей заключаются в следующих утверждениях.



\begin{enumerate}

  \item Неособая проективная поверхность $X$ тогда и только тогда является относительно М. М., когда она не содержит исілючительных кривых первого рода (см. Исключителъное подмногообразие).



  \item Всякая неособая полная поверхность обладает бирациональным морфизмом на относительно М. м.



  \item В каждом непустом классе $B$ бирационально әквивалентных поверхностей, кроме классов рациональных и линейтатых поверхностей, существует (и притом единственная) М. м.



  \item Если $B$ - класс линейчатых поверхностей с кривой $C$ рода $g>0$ в качестве базы, то все относительно M. м. в $B$ исчерпываются геометрически линейчатыми поверхностями $\pi: X \rightarrow C$.



  \item Если $B$ - класс рациональных поверхностей, то все относительно М. м. в $B$ исчерпывалтся проективной плоскостью $P^{2}$ и серией минимальных рациональных линейчатых поверхностей $F_{n}=P\left(\widehat{G}_{p 1}+\widehat{G}_{p 1}(n)\right)$ для всех целых $n \geqslant 2$.



\end{enumerate}



Имеется (см. [6], [7]) обобщение теории М. м. поверхностей на регулярные двумерные схемы. Описаны (см. [2]) М. м. рациональных поверхностей, определенных над произвольным полем.



О М. м. Для многообразий размерности $n \geqslant 3$ почти ничего не известно (1982).



Лит.: [1] Алгебраические поверхности, М., 1965 (Тр. Матем. ин-та АН ССCP, т. 75); [2] И с к о в с к и х B. А., "Изв. АН CCCP. Cep. матем.», 1979, т. 43, No 1, c. 19-43; [3] III a \begin{CJK}{UTF8}{mj}中\end{CJK} a p e-



\includegraphics[max width=\textwidth, center]{2023_07_08_f09922783b3f583efe9cg-1(3)}

embeddings of surfaces. Proc. symp. pure math., $\mathrm{v}$. 29 . Algebraic



\begin{center}

\includegraphics[max width=\textwidth]{2023_07_08_f09922783b3f583efe9cg-1(2)}

\end{center}



\includegraphics[max width=\textwidth, center]{2023_07_08_f09922783b3f583efe9cg-1(1)}

p. $380-405$; [7] $\mathrm{S} \mathrm{h}$ a f a r e v i t $\mathrm{ch}$ I., Lectures on minimal models and birational transformations of twodimensional schemes, MDey 1966. МИНИМАЛЬНАЯ НОРМАЛЬНАЯ ПОДГРУ ППА- -

пеединичная нормальная подгруппа $H$ такая, что между ней и единичной подгруппой нет других нормальных подгрупп всей группы. М. н. п. имеются далеко не во всякой группе. Если группа конечна, то любая ее М. н. п. является прямым произведением изоморфных простых групп. Если М. н. П. у группы $G$ существует и единственна, то она наз. м о н о ли т о м (иногда с ердце виной), а сама группа $G$ наз. монолит и ч н о ї.



Лит.: [1] К у р о ші А. Г., Теория групп, 3 изд., М., 1967.



МИНИМАЛЬНАЯ ПОВЕРХНОСТЬ - Поверхность, у к-рой средняя кривизна $H$ равна нулю во всех точках.



Первые исследования о М. п. восходят к Ж. Лагранжу (J. Lagrange, 1768), к-рый рассмотрел следующую вариационную задачу: найти поверхность наименьшей площади, натянутую на данный контур. Предполагая искомую поверхность задаваемой в виде $z=z(x, y)$, Ж. Јагранж получил, что $z(x, y)$ необходимо должна удовлетворять т. н. у р а н е н и ю Э й ле р аЛ а г р а н ж а



\[

\begin{gathered}

\left(1+q^{2}\right) \frac{\partial^{2} z}{\partial x^{2}}-2 p q \frac{\partial^{2} z}{\partial x \partial y}+\left(1+p^{2}\right) \frac{\partial^{2} z}{\partial y^{2}}=0, \\

p=\frac{\partial z}{\partial x}, \quad q=\frac{\partial z}{\partial y} .

\end{gathered}

\]





\end{document}